\chapter{PLANOS Y LAMINAS DE DETALLE}

Los planos deben colocarse en la carpeta:

\begin{center}
\texttt{latex/planos/}
\end{center}

Cuando tengas cada plano en PDF o imagen, guarda el archivo con el nombre indicado y actualiza el estado en la Tabla~\ref{tab:planos}.

\begin{longtable}{p{1.6cm}p{4cm}p{5cm}p{2cm}}
\caption{Relacion de planos requeridos}\label{tab:planos}\\
\toprule
\textbf{Codigo} & \textbf{Plano} & \textbf{Archivo esperado} & \textbf{Estado} \\
\midrule
\endfirsthead
\toprule
\textbf{Codigo} & \textbf{Plano} & \textbf{Archivo esperado} & \textbf{Estado} \\
\midrule
\endhead
IE-01 & Ubicacion y arquitectura/croquis & \texttt{IE-01-ubicacion-arquitectura.pdf} & Pendiente \\
IE-02 & Alumbrado & \texttt{IE-02-alumbrado.pdf} & Pendiente \\
IE-03 & Tomacorrientes & \texttt{IE-03-tomacorrientes.pdf} & Pendiente \\
IE-04 & Circuitos y canalizaciones & \texttt{IE-04-circuitos-canalizaciones.pdf} & Pendiente \\
IE-05 & Diagrama unifilar y cuadro de cargas & \texttt{IE-05-unifilar-cuadro-cargas.pdf} & Pendiente \\
IE-06 & Puesta a tierra & \texttt{IE-06-puesta-tierra.pdf} & Pendiente \\
\bottomrule
\end{longtable}

\section{Planos a insertar}

\subsection{IE-01: Ubicacion y arquitectura}

\editar{Insertar aqui el plano o croquis de ubicacion y arquitectura.}

% Cuando tengas el plano, descomenta y ajusta:
% \begin{figure}[H]
% \centering
% \includegraphics[width=\textwidth]{planos/IE-01-ubicacion-arquitectura.pdf}
% \caption{Plano de ubicacion y arquitectura}
% \end{figure}

\subsection{IE-02: Alumbrado}

\editar{Insertar aqui el plano de alumbrado.}

\subsection{IE-03: Tomacorrientes}

\editar{Insertar aqui el plano de tomacorrientes.}

\subsection{IE-04: Circuitos y canalizaciones}

\editar{Insertar aqui el plano de circuitos y canalizaciones.}

\subsection{IE-05: Diagrama unifilar y cuadro de cargas}

\editar{Insertar aqui el diagrama unifilar del tablero general y el cuadro de cargas.}

\subsection{IE-06: Puesta a tierra}

\editar{Insertar aqui el detalle de puesta a tierra.}

\section{Detalles sugeridos}

\begin{itemize}
  \item Simbologia electrica usada.
  \item Leyenda de circuitos.
  \item Altura de interruptores y tomacorrientes, si el docente lo pide.
  \item Detalle del tablero.
  \item Detalle de puesta a tierra.
\end{itemize}

